

\illustration{0.3}{watch.jpg}{
    CC BY-SA 4.0 by Jordiferrer on \href{https://commons.wikimedia.org/wiki/File:H4_timekeeper.jpg}{Wikimedia Commons}
}

% optionally define variables/limits for this problem

% See https://en.wikipedia.org/wiki/History_of_timekeeping_devices#Development_of_the_watch
It is the 25th of November, 1625.
You have put a lot of time into making a mechanical watch,
but you are unsure if the hands are placed correctly.
Even a watch that is not running is correct twice a day,
but if the hands of the watch do not correspond to a real time,
it is never correct at all.
Without wasting any more time, you measure the angles of the hands
and check whether they correspond to a real time.

For example, consider the first sample input, visualized in Figure~\ref{fig:clocks}:
if the minute hand points to a quarter past (90~degrees from 12~o'clock),
then the hour hand cannot point exactly to 2~o'clock (60~degrees from 12~o'clock).

You know that both hands of the watch are moving continuously.

\begin{figure}[h]
    \centering
    \hfill
    \includegraphics[width=0.2\textwidth]{clock1.pdf}
    \hfill
    \includegraphics[width=0.2\textwidth]{clock2.pdf}
    \hfill
    \includegraphics[width=0.2\textwidth]{clock3.pdf}
    \hfill
    \includegraphics[width=0.2\textwidth]{clock4.pdf}
    \hfill
    \hfill
    \caption{
        Illustrations of the sample inputs.
        The first and third sample inputs \emph{do not} correspond to a real time.
        The second and fourth sample inputs \emph{do} correspond to a real time.
    }
    \label{fig:clocks}
\end{figure}

